\documentclass{ibetest}
\begin{document}
\maketest{Progress Test 1.A --- Thermo-S1}
         {1.A \,$\cdot$\, System and state variables}
         {7 minutes}{14}

% ---------------------------------------------------------------------------
\exercise{Unit conversion}{1 mark}
Convert the molar volume $V_m = \qty{24.5}{\litre\per\mole}$ into
$\mathrm{m^3\,mol^{-1}}$. Give the exact result, without rounding.
\answer{1}{2.2cm}{%
  $1\ \mathrm L = 1\ \mathrm{dm^3} = (10^{-1}\ \mathrm m)^3 = 10^{-3}\ \mathrm{m^3}$
  \qquad $V_m = 24.5\times10^{-3}\ \mathrm{m^3\,mol^{-1}}$ \qquad
  $\boxed{V_m = 2.45\times10^{-2}\ \mathrm{m^3\,mol^{-1}}}$}
\begin{scheme}
\pt{1}{1 pt}{$2.45\times10^{-2}\ \mathrm{m^3\,mol^{-1}}$, with the unit written correctly.}
\end{scheme}
\begin{tip}
A litre is a cubic decimetre: cube the conversion factor, $(10^{-1})^3 = 10^{-3}$. Every
SI formula of this module ($pV = nRT$, \dots) needs volumes in $\mathrm{m^3}$.
\end{tip}

% ---------------------------------------------------------------------------
\exercise{Closed and isolated systems --- definitions}{5 marks}
Define a closed system and an isolated system. For a gas enclosed in a cylinder, say what the
walls must be like to prevent each kind of exchange with the surroundings. Is every closed
system isolated?
\answer{2,2,1}{6.0cm}{%
  \textbf{Closed system:} it exchanges no \emph{matter} with its surroundings
  ($n$ is constant).\\
  \textbf{Isolated system:} it exchanges neither \emph{matter} nor \emph{energy} with its
  surroundings.\\[4pt]
  Walls needed for the gas to be isolated:
  \begin{itemize}
  \item \textbf{sealed} walls, so that no matter is exchanged;
  \item \textbf{rigid} walls and a \textbf{fixed} piston, so that no work is exchanged
        (pressure forces cannot move them);
  \item \textbf{thermally insulating} walls, so that no heat is exchanged.
  \end{itemize}
  An isolated system is necessarily closed, but \textbf{the converse is false}. For example,
  a sealed weather balloon is closed, yet it exchanges heat and work with the air around it.}
\begin{scheme}
\pt{1}{2 pts}{Closed: no exchange of matter (1~pt). Isolated: no matter and no energy (1~pt).}
\pt{2}{2 pts}{Sealed, rigid or fixed, and insulating walls (1~pt if only two of the three).}
\pt{3}{1 pt}{Isolated $\Rightarrow$ closed, but closed $\not\Rightarrow$ isolated, with an
  example.}
\end{scheme}

% ---------------------------------------------------------------------------
\exercise{Classifying quantities and systems}{4 marks}
\question{Which of these quantities is intensive?}
\opttwo{0}{the mass $m$}{0}{the volume $V$}{1}{the molar volume $V_m = V/n$}{0}{the energy $E$}
\why{the quotient of two extensive quantities is intensive: doubling the system doubles
  both $V$ and $n$, so $V/n$ does not change.}
\question{A hot-air balloon, with its opening at the bottom, is:}
\opts{0}{closed}{1}{open}{0}{isolated}{0}{closed and isolated}
\why{air enters and leaves through the opening, so the system exchanges matter
  (Fig.~1 of the notes).}
\question{The SI dimension of pressure is:}
\opts{1}{$\mathrm{M\,L^{-1}\,T^{-2}}$}{0}{$\mathrm{M\,L\,T^{-2}}$}{0}{$\mathrm{M\,L^{2}\,T^{-2}}$}
     {0}{$\mathrm{M\,L^{-2}\,T^{-2}}$}
\why{pressure is a force per unit area:
  $\mathrm{M\,L\,T^{-2}}/\mathrm{L^2} = \mathrm{M\,L^{-1}\,T^{-2}}$ (Tab.~1 of the notes).}
\question{$R$ is an intensive constant. Which relation is homogeneous in the extensive /
  intensive sense?}
\opts{0}{$\dfrac pT = \dfrac RV$}{0}{$pV = RT$}{1}{$\dfrac pT = \dfrac{nR}{V}$}{0}{$p = nRT$}
\why{$p/T$ is intensive. So is $nR/V$ (an extensive quantity divided by an extensive one).
  In the other options, one side doubles with the system and the other does not (Box~7).}

% ---------------------------------------------------------------------------
\exercise{Composition of a solution}{4 marks}
A mass $m$ of sodium chloride (molar mass $M$) is dissolved in water to obtain a volume $V$
of solution. Compute the amount of substance $n$ of solute, the mass concentration $c_m$ and
the molar concentration $c$. Check that $c_m = cM$.
\data{$m = \qty{5.85}{g}$ \hspace{14mm} $M = \qty{58.5}{\gram\per\mole}$ \hspace{14mm}
      $V = \qty{0.500}{L}$}
\answer{1,1,1,1}{5.6cm}{%
  $n = \dfrac mM = \dfrac{5.85\ \mathrm g}{58.5\ \mathrm{g\,mol^{-1}}}$ \qquad
  $\boxed{n = 0.100\ \mathrm{mol}}$\\[4pt]
  $c_m = \dfrac mV = \dfrac{5.85\ \mathrm g}{0.500\ \mathrm L}$ \qquad
  $\boxed{c_m = 11.7\ \mathrm{g\,L^{-1}}}$\\[4pt]
  $c = \dfrac nV = \dfrac{0.100\ \mathrm{mol}}{0.500\ \mathrm L}$ \qquad
  $\boxed{c = 0.200\ \mathrm{mol\,L^{-1}}}$\\[4pt]
  Check: $cM = 0.200\times58.5 = 11.7\ \mathrm{g\,L^{-1}} = c_m$. \checkmark\ Both
  concentrations are intensive (quotients of extensive quantities).}
\begin{scheme}
\pt{1}{1 pt}{$n = m/M = 0.100\ \mathrm{mol}$.}
\pt{2}{1 pt}{$c_m = 11.7\ \mathrm{g\,L^{-1}}$.}
\pt{3}{1 pt}{$c = 0.200\ \mathrm{mol\,L^{-1}}$.}
\pt{4}{1 pt}{$c_m = cM$ checked.}
\end{scheme}
\begin{tip}
Test~1, Ex.~4: write $n = m/M$ with letters first, then substitute \emph{with units}. The
units then cancel correctly by themselves: $\mathrm g/(\mathrm{g\,mol^{-1}}) = \mathrm{mol}$.
\end{tip}

\finish
\end{document}
